\documentclass{article}
\usepackage[UTF8]{ctex}
\usepackage[tc]{titlepic}
\usepackage{titlesec}
\usepackage{cite}
\usepackage{fancyhdr}
\usepackage{booktabs}
\usepackage{graphicx}
\usepackage{geometry}
\usepackage[section]{placeins}
\usepackage{amsmath,amssymb}
\usepackage{longtable}
\usepackage{enumitem}
\usepackage{float}
\geometry{a4paper,scale=0.8}
\pagestyle{fancy}

\lhead{数学建模课程\\\today}
\chead{中国科学技术大学\\数学建模课程}
\rhead{Predator--Prey Fear Effect\\ {\CTEXoptions[today=old]\today}}
\newcommand{\upcite}[1]{\textsuperscript{\cite{#1}}}
\graphicspath{{figures/}{../figures/}}

\titleformat*{\section}{\bfseries\Large}
\titleformat*{\subsection}{\bfseries\large}
\titleformat*{\subsubsection}{\bfseries\normalsize}

\title{\bfseries 恐惧效应与记忆反馈下捕食者--猎物系统的动力学建模与稳定性分析}
\author{姓名：\underline{\hspace{3cm}} \quad 课程ID：\underline{\hspace{2.5cm}} \quad 学号：\underline{\hspace{3cm}}}
\date{\today}

\begin{document}
\maketitle

\begin{abstract}
本文围绕恐惧效应下的捕食者--猎物动力学建立机制模型。模型主线为：捕食者数量改变猎物感知风险，感知风险降低猎物有效出生率，并进一步改变猎物与捕食者的长期动力学。本文在经典 Holling II 捕食者--猎物模型基础上，依次建立无恐惧模型、即时恐惧模型和记忆恐惧模型，并将三级食物链作为扩展模块。理论部分证明正性和有界性，推导正共存平衡、Jacobian、即时恐惧模型局部稳定性以及记忆模型的 Routh--Hurwitz 条件。数值部分通过时间序列、相图、$k$ 分岔图、$k-\alpha$ 二参数热力图和理论边界对照，展示恐惧稳定系统、恐惧降低捕食者密度但仍保持振荡、长记忆诱发或放大振荡三类典型现象。结果说明：恐惧效应不应简单理解为“使猎物更少”或“必然稳定系统”，其作用方向依赖参数区间和记忆反馈速度；本文结论属于机制模拟和定性分析，不声称对真实生态系统进行精确预测。
\end{abstract}

\clearpage
\setcounter{section}{1}
\section*{\centerline{一、前言（问题的提出）}}

捕食者对猎物的影响通常首先被理解为直接捕食，即捕食者通过杀死并消耗猎物减少猎物数量。但生态学研究表明，捕食者还会通过非消耗效应影响猎物：猎物感知到捕食风险后，可能减少觅食、改变活动区域、提高警觉性，甚至降低繁殖投入。Lima 与 Dill\upcite{lima1990behavioral}指出，动物会在生态时间尺度上评估和调节自身面临的捕食风险。Preisser 与 Bolnick\upcite{preisser2008many}进一步将这类作用概括为非消耗性捕食者效应。Zanette 等\upcite{zanette2011perceived}的鸣禽实验显示，仅感知捕食风险就可以使年后代数量下降约 $40\%$。这说明捕食者即使没有直接捕食，也可能显著改变猎物种群增长。

因此，本项目关注如下机制链条：
\[
\text{捕食者数量}
\longrightarrow
\text{猎物感知风险}
\longrightarrow
\text{猎物繁殖下降}
\longrightarrow
\text{种群动力学改变}.
\]
其中，感知风险不仅可以由当前捕食者数量决定，还可能受到历史捕食压力影响。为避免模型过度复杂，本文先完成二维捕食者--猎物系统中的恐惧效应和记忆反馈建模，再将三级食物链作为扩展讨论。

\setcounter{section}{2}
\section*{\centerline{二、相关工作}}

\subsection{行为生态学依据}
恐惧效应的生态学基础来自行为生态学。Lima 与 Dill\upcite{lima1990behavioral}强调动物会在觅食、逃避、社群活动和繁殖行为中权衡捕食风险与收益。Preisser 与 Bolnick\upcite{preisser2008many}指出，非消耗效应可通过行为、生理、存活和繁殖等路径影响猎物种群。Zanette 等\upcite{zanette2011perceived}通过野外实验说明，在排除直接捕食后，仅捕食者声音造成的风险感知即可显著降低鸣禽繁殖成功率。本文用该结果对恐惧函数的量级进行说明，而不进行真实参数拟合。

\subsection{经典恐惧模型}
Wang、Zanette 与 Zou\upcite{wang2016modelling}是本文主要数学参考。该文将恐惧函数 $f(k,y)$ 乘入猎物出生项，常用形式为
\[
f(k,y)=\frac{1}{1+ky}.
\]
在 Holling II 捕食函数下，恐惧强度不仅会改变共存平衡稳定性，还可能改变 Hopf 分岔方向；高恐惧可稳定系统，低恐惧则可能产生周期解、多重极限环或双稳态。

\subsection{时滞、记忆与扩展模型}
Liu 等\upcite{liu2021dynamics}在恐惧效应基础上加入妊娠时滞，给出了可借鉴的参数组和 $k$ 分岔图。Yang 与 Jin\upcite{yang2022dynamics}研究的是带空间记忆的扩散型捕食者--猎物模型，本文不复现其 PDE，而是将“历史风险影响当前动力学”的思想简化为 ODE 记忆变量。Panday 等\upcite{panday2018stability}研究了三级食物链中的恐惧效应，说明恐惧可能使复杂食物链动力学趋于稳定。Zhao、Yu 与 Li\upcite{zhao2025dynamics}提供了空间扩散和 Turing 不稳定的近期扩展参考。

\setcounter{section}{3}
\section*{\centerline{三、问题分析}}

题目要求从机制出发建立模型，而不是用真实数据做精确预测。因此，本文将主要问题拆成三层：
\begin{enumerate}[label=(\arabic*)]
    \item 无恐惧时，Holling II 捕食者--猎物模型是否产生稳定共存或周期振荡；
    \item 引入即时恐惧后，恐惧强度 $k$ 如何改变平衡点、捕食者密度和振荡幅度；
    \item 引入记忆变量 $z(t)$ 后，记忆更新速率 $\alpha$ 如何改变稳定性和周期行为。
\end{enumerate}

本文重点不是把所有机制都塞入同一个主模型，而是保持主模型足够清晰：恐惧只降低猎物有效出生率，捕食仍通过 Holling II 函数发生，记忆变量只表示风险感知，不代表第三个物种。三级食物链扩展只用于说明多营养级恐惧效应的新问题。

\setcounter{section}{4}
\section*{\centerline{四、建模的假设}}

\begin{enumerate}[label=\textbf{H\arabic*.}]
    \item 系统封闭，不考虑外来迁入和迁出。
    \item 空间环境均匀，不考虑空间扩散和斑图形成。
    \item 时间连续，种群密度用连续变量表示。
    \item 猎物在无捕食者时具有出生、自然死亡和种内竞争。
    \item 捕食者依赖猎物获得增长，死亡率为常数。
    \item 捕食函数采用 Holling II 型。
    \item 恐惧效应不直接杀死猎物，而是降低猎物有效出生率。
    \item 猎物对捕食风险有记忆，记忆风险 $z(t)$ 是历史捕食者密度的平滑平均。
    \item 所有参数非负，主要分析正参数情况。
    \item 初始种群密度为正。
    \item 数值模拟参考文献参数和合理区间，不进行真实参数拟合。
    \item 若使用经验数据，只用于现象说明或参数量级校准，不声称精确预测真实种群。
\end{enumerate}

\setcounter{section}{5}
\section*{\centerline{五、符号说明}}

\begin{table}[H]
    \caption{\textbf{主模型符号说明}}
    \centering
    \begin{tabular}{ccc}
    \toprule
    符号 & 说明 & 单位或属性 \\
    \midrule
    $x(t)$ & 猎物密度 & 连续变量 \\
    $y(t)$ & 捕食者密度 & 连续变量 \\
    $z(t)$ & 猎物感知到的捕食风险/记忆风险 & 连续变量 \\
    $r$ & 猎物最大出生率 & 正参数 \\
    $d_1$ & 猎物自然死亡率 & 非负参数 \\
    $d_2$ & 猎物种内竞争强度 & 正参数 \\
    $p$ & 捕食强度 & 正参数 \\
    $h$ & Holling II 处理时间参数 & 非负参数 \\
    $\eta$ & 捕食转化效率 & 正参数 \\
    $d_3$ & 捕食者死亡率 & 正参数 \\
    $k$ & 恐惧强度 & 非负参数 \\
    $\alpha$ & 记忆更新速率，越大表示记忆越短 & 正参数 \\
    \bottomrule
    \end{tabular}
\end{table}

\setcounter{section}{6}
\section*{\centerline{六、数学模型建立与理论分析}}

\subsection{无恐惧基准模型 M0}
令 $k=0$，不引入记忆变量，得到
\begin{equation}
\begin{cases}
\dfrac{dx}{dt}=rx-d_1x-d_2x^2-\dfrac{pxy}{1+hx},\\[6pt]
\dfrac{dy}{dt}=\eta\dfrac{pxy}{1+hx}-d_3y.
\end{cases}
\end{equation}
该模型作为后续恐惧模型的对照。

\subsection{即时恐惧模型 M1}
即时恐惧假设猎物感知风险等于当前捕食者密度，即 $z=y$：
\begin{equation}
\begin{cases}
\dfrac{dx}{dt}=\dfrac{rx}{1+ky}-d_1x-d_2x^2-\dfrac{pxy}{1+hx},\\[6pt]
\dfrac{dy}{dt}=\eta\dfrac{pxy}{1+hx}-d_3y.
\end{cases}
\end{equation}
其中 $1/(1+ky)$ 表示捕食风险降低猎物有效出生率。

\subsection{记忆恐惧模型 M2}
为描述历史风险，本文引入记忆变量 $z(t)$：
\begin{equation}
\label{eq:m2}
\begin{cases}
\dfrac{dx}{dt}=\dfrac{rx}{1+kz}-d_1x-d_2x^2-\dfrac{pxy}{1+hx},\\[6pt]
\dfrac{dy}{dt}=\eta\dfrac{pxy}{1+hx}-d_3y,\\[6pt]
\dfrac{dz}{dt}=\alpha(y-z).
\end{cases}
\end{equation}
当 $\alpha$ 较大时，$z$ 快速跟随 $y$，记忆较短；当 $\alpha$ 较小时，风险感知更新慢，历史压力的影响更长。

\subsection{正性和有界性}
若初始值 $x(0)>0,y(0)>0,z(0)>0$，则在坐标平面 $x=0$ 和 $y=0$ 上均有对应导数为零，轨道不会穿出非负区域；当 $z=0$ 且 $y>0$ 时，$z'=\alpha y>0$。因此正初值保持正性。

又因为恐惧函数不超过 $1$，且捕食项非负，
\[
\frac{dx}{dt}\le (r-d_1)x-d_2x^2.
\]
当 $r>d_1$ 时，
\[
\limsup_{t\to\infty}x(t)\le \frac{r-d_1}{d_2}.
\]
令 $W(t)=\eta x(t)+y(t)$，直接捕食项在组合函数中抵消，可得到 $W$ 受线性耗散不等式控制，因此 $y(t)$ 最终有界。由于 $z'=\alpha(y-z)$ 是对 $y$ 的指数加权平均，$y$ 有界则 $z$ 也有界。

\subsection{平衡点}
灭绝平衡为 $E_0=(0,0,0)$。若 $r>d_1$，猎物单独存在平衡为
\[
E_1=\left(\frac{r-d_1}{d_2},0,0\right).
\]
正平衡满足 $z^\ast=y^\ast$。由捕食者方程可得
\[
\eta\frac{p x^\ast}{1+h x^\ast}=d_3.
\]
若 $\eta p>d_3h$，则
\[
x^\ast=\frac{d_3}{\eta p-d_3h}.
\]
令
\[
B=\frac{p}{1+h x^\ast},\qquad H=d_1+d_2x^\ast,
\]
猎物方程化为
\[
Bk(y^\ast)^2+(B+kH)y^\ast+(H-r)=0.
\]
当 $r>H$ 时存在正根
\[
y^\ast=
\frac{-(B+kH)+\sqrt{(B+kH)^2+4Bk(r-H)}}{2Bk},
\qquad z^\ast=y^\ast.
\]
当 $k=0$ 时，退化为 $y^\ast=(r-H)/B$。

\subsection{Jacobian 与稳定性}
M1 在正平衡处的 Jacobian 为
\[
J=
\begin{pmatrix}
F_x & F_y\\
G_x & G_y
\end{pmatrix},
\]
其中
\[
F_x=\frac{r}{1+ky^\ast}-d_1-2d_2x^\ast-\frac{p y^\ast}{(1+h x^\ast)^2},
\]
\[
F_y=-\frac{rkx^\ast}{(1+ky^\ast)^2}-\frac{px^\ast}{1+h x^\ast}<0,
\]
\[
G_x=\frac{\eta p y^\ast}{(1+h x^\ast)^2}>0,\qquad
G_y=\eta\frac{p x^\ast}{1+hx^\ast}-d_3=0.
\]
因此 $\det J=-F_yG_x>0$，局部稳定性主要由 $\operatorname{tr}J=F_x$ 决定。若 $F_x<0$，正平衡局部渐近稳定；若 $F_x>0$，正平衡不稳定；若 $F_x=0$，可能发生 Hopf 分岔。

M2 在正平衡处的 Jacobian 可写成
\[
J=
\begin{pmatrix}
A & B & C\\
D & 0 & 0\\
0 & \alpha & -\alpha
\end{pmatrix},
\]
其中
\[
A=\frac{r}{1+kz^\ast}-d_1-2d_2x^\ast-\frac{p y^\ast}{(1+h x^\ast)^2},
\quad
B=-\frac{px^\ast}{1+hx^\ast},
\]
\[
C=-\frac{rkx^\ast}{(1+kz^\ast)^2},
\quad
D=\frac{\eta p y^\ast}{(1+h x^\ast)^2}.
\]
其特征多项式为
\[
\lambda^3+A_1\lambda^2+A_2\lambda+A_3=0,
\]
其中
\[
A_1=\alpha-A,\qquad
A_2=-A\alpha-BD,\qquad
A_3=-\alpha D(B+C).
\]
Routh--Hurwitz 稳定条件为
\[
A_1>0,\qquad A_2>0,\qquad A_3>0,\qquad A_1A_2>A_3.
\]
Hopf 临界边界可由 $A_1A_2=A_3$ 给出。

\begin{figure}[H]
    \centering
    \includegraphics[width=0.82\textwidth]{fig01_mechanism_diagram.png}
    \caption{恐惧效应与直接捕食路径的机制示意图}
    \label{fig:mechanism}
\end{figure}

\begin{figure}[H]
    \centering
    \includegraphics[width=0.78\textwidth]{fig02_model_layers.png}
    \caption{本文使用的模型层级：M0、M1、M2 与 M3}
    \label{fig:layers}
\end{figure}

\setcounter{section}{7}
\section*{\centerline{七、结果（与对比）}}

\subsection{模拟设置}
数值求解优先使用 \texttt{scipy.integrate.solve\_ivp} 的 DOP853 方法，容差设为 \texttt{rtol=1e-9, atol=1e-11}；二维大网格扫描采用向量化 RK4 以提高复现效率。尾部区间计算均值、最大值、最小值、振幅和相对振幅，并按照稳定共存、振荡共存、捕食者灭绝、猎物灭绝和低密度风险进行分类。

主要参数组来自当前已有强振荡参数组的文献型转换：
\[
r=0.5,\quad d_1=0,\quad d_2=0.005,\quad p=0.08,\quad h=0.1,\quad \eta=0.5,\quad d_3=0.3.
\]
这里 $p=0.08,h=0.1$ 对应原 Rosenzweig--MacArthur 形式 $0.8xy/(10+x)$。

\begin{table}[H]
    \caption{\textbf{主数值实验参数}}
    \centering
    \begin{tabular}{cccccccccc}
    \toprule
    $r$ & $d_1$ & $d_2$ & $p$ & $h$ & $\eta$ & $d_3$ & $x_0$ & $y_0$ & $z_0$ \\
    \midrule
    0.5 & 0 & 0.005 & 0.08 & 0.1 & 0.5 & 0.3 & 30 & 10 & 10 \\
    \bottomrule
    \end{tabular}
    \label{tab:params}
\end{table}

\subsection{时间序列与相图}
图 \ref{fig:timeseries} 对比了无恐惧、即时恐惧、快记忆恐惧和慢记忆恐惧下的时间序列。无恐惧时系统产生明显周期振荡；加入恐惧后，捕食者峰值和振荡幅度明显改变。慢记忆情况下，风险感知滞后于当前捕食者数量，可能导致振荡保留甚至被放大。

\begin{figure}[H]
    \centering
    \includegraphics[width=\textwidth]{fig03_time_series_cases.png}
    \caption{不同模型和代表性参数下的时间序列}
    \label{fig:timeseries}
\end{figure}

\begin{figure}[H]
    \centering
    \includegraphics[width=0.95\textwidth]{fig04_phase_portraits.png}
    \caption{尾部区间的猎物--捕食者相图}
    \label{fig:phase}
\end{figure}

\subsection{$k$ 分岔图与均值振幅}
图 \ref{fig:bifurcation} 展示了尾部最大值和最小值随恐惧强度 $k$ 的变化。当最大值和最小值分离时，说明系统仍有周期振荡；当两者合并时，系统趋于稳定平衡。图 \ref{fig:meanamp} 显示，在本参数组下，$k$ 增大通常降低捕食者平均密度和振荡幅度；无恐惧时捕食者尾部振幅较大，而高恐惧下振幅趋近于零。

\begin{table}[H]
    \caption{\textbf{$k$ 扫描中的代表点}}
    \centering
    \begin{tabular}{ccccc}
    \toprule
    $k$ & 分类 & $\bar{x}$ & $\bar{y}$ & $A_y$ \\
    \midrule
    0.0 & 振荡共存 & 58.727 & 9.563 & 25.393 \\
    0.2 & 稳定共存 & 30.000 & 5.000 & $1.69\times10^{-10}$ \\
    1.0 & 稳定共存 & 30.000 & 1.713 & $1.08\times10^{-11}$ \\
    5.2 & 稳定共存 & 30.000 & 0.415 & $2.48\times10^{-12}$ \\
    20.0 & 稳定共存 & 30.000 & 0.114 & $7.08\times10^{-13}$ \\
    \bottomrule
    \end{tabular}
    \label{tab:scan-k}
\end{table}

\begin{figure}[H]
    \centering
    \includegraphics[width=0.86\textwidth]{fig05_bifurcation_k_extrema.png}
    \caption{恐惧强度 $k$ 的分岔式尾部极值图}
    \label{fig:bifurcation}
\end{figure}

\begin{figure}[H]
    \centering
    \includegraphics[width=0.86\textwidth]{fig06_mean_amplitude_vs_k.png}
    \caption{均值和振幅随恐惧强度 $k$ 的变化}
    \label{fig:meanamp}
\end{figure}

\subsection{$k-\alpha$ 二参数扫描}
图 \ref{fig:classheat} 显示了 $k-\alpha$ 参数平面中的动力学分类。扫描结果包含稳定共存、振荡共存、低密度风险和少量捕食者灭绝区域。图 \ref{fig:ampheat} 用连续振幅指标展示同一参数平面，可以看出不同振荡区域的强弱差异。

\begin{table}[H]
    \caption{\textbf{$k-\alpha$ 网格扫描分类计数}}
    \centering
    \begin{tabular}{cc}
    \toprule
    分类 & 网格点数量 \\
    \midrule
    稳定共存 stable\_coexistence & 4550 \\
    振荡共存 oscillatory\_coexistence & 850 \\
    低密度风险 low\_density\_risk & 156 \\
    捕食者灭绝 predator\_extinct & 9 \\
    数值无效 invalid & 2616 \\
    \bottomrule
    \end{tabular}
    \label{tab:class-count}
\end{table}

\begin{figure}[H]
    \centering
    \includegraphics[width=0.88\textwidth]{fig07_k_alpha_classification_heatmap.png}
    \caption{$k-\alpha$ 参数平面的动力学分类热力图}
    \label{fig:classheat}
\end{figure}

\begin{figure}[H]
    \centering
    \includegraphics[width=0.88\textwidth]{fig08_k_alpha_amplitude_heatmap.png}
    \caption{$k-\alpha$ 参数平面的连续振幅热力图}
    \label{fig:ampheat}
\end{figure}

\subsection{理论边界与数值结果对比}
图 \ref{fig:theory} 将 Routh--Hurwitz 条件给出的临界边界 $A_1A_2=A_3$ 叠加到数值分类热力图上。理论边界能够解释稳定区和振荡区的大体分界；在低密度风险、长周期振荡和有限模拟时间附近，理论线性边界与数值分类会出现局部差异，这是非线性系统数值分类中的常见现象。

\begin{figure}[H]
    \centering
    \includegraphics[width=0.88\textwidth]{fig09_theory_boundary_vs_simulation.png}
    \caption{Routh--Hurwitz 理论边界与数值分类对照}
    \label{fig:theory}
\end{figure}

\subsection{恐惧函数鲁棒性与经验量级校准}
为检验结论是否依赖单一恐惧函数，本文比较三种写法：
\[
f_1(k,z)=\frac{1}{1+kz},\qquad
f_2(k,z)=e^{-kz},\qquad
f_3(k,z)=\frac{1}{1+kz+k_2z^2}.
\]
图 \ref{fig:robust} 表明，替换恐惧函数后，捕食者平均密度和振幅随 $k$ 变化的总体趋势仍然保留，说明主要结论具有一定稳健性。

Zanette 等\upcite{zanette2011perceived}报告的 $40\%$ 繁殖下降可作为恐惧函数量级参考。若将高风险归一化为 $z=1$，并要求
\[
\frac{1}{1+k}=0.6,
\]
则 $k\approx 0.667$。图 \ref{fig:calibration} 只用于说明量级，不用于真实拟合。

\begin{figure}[H]
    \centering
    \includegraphics[width=0.86\textwidth]{fig10_fear_function_robustness.png}
    \caption{不同恐惧函数形式下的结果鲁棒性}
    \label{fig:robust}
\end{figure}

\begin{figure}[H]
    \centering
    \includegraphics[width=0.78\textwidth]{fig11_dataset_or_empirical_calibration.png}
    \caption{基于 Zanette 等实验结果的恐惧函数量级校准}
    \label{fig:calibration}
\end{figure}

\subsection{代表性现象}
参数扫描自动筛选出三类代表性行为：
\begin{itemize}
    \item Case A：无恐惧时系统振荡，较强恐惧使系统转为稳定共存；
    \item Case B：慢记忆下恐惧降低捕食者峰值，但系统仍保持振荡；
    \item Case C：长记忆诱发或放大振荡，而快记忆条件下系统稳定。
\end{itemize}
这三类结果说明，恐惧效应的作用并非单调地“稳定”或“破坏”系统，而取决于恐惧强度与记忆更新速度共同决定的反馈结构。

\begin{table}[H]
    \caption{\textbf{自动搜索得到的代表性行为}}
    \centering
    \begin{tabular}{cccccc}
    \toprule
    案例 & $k$ & $\alpha$ & 分类 & $\bar{y}$ & $A_y$ \\
    \midrule
    Case A & 1.0 & 5.012 & 稳定共存 & 1.713 & 0.000 \\
    Case B & 4.8 & 0.010 & 振荡共存 & 0.421 & 0.517 \\
    Case C & 5.2 & 0.010 & 振荡共存 & 0.428 & 1.502 \\
    \bottomrule
    \end{tabular}
    \label{tab:cases}
\end{table}

\subsection{三级食物链扩展}
三级食物链扩展模型为
\begin{equation}
\begin{cases}
\dfrac{dX}{dt}=
\dfrac{rX}{1+k_1Y}
-d_X X-a_X X^2
-\dfrac{p_1XY}{1+h_1X},\\[6pt]
\dfrac{dY}{dt}=
\dfrac{\eta_1p_1XY}{1+h_1X}
\cdot
\dfrac{1}{1+k_2Z}
-d_YY
-\dfrac{p_2YZ}{1+h_2Y},\\[6pt]
\dfrac{dZ}{dt}=
\eta_2\dfrac{p_2YZ}{1+h_2Y}
-d_ZZ.
\end{cases}
\end{equation}
其中 $k_1$ 表示底层猎物对中间捕食者的恐惧，$k_2$ 表示中间捕食者对顶级捕食者的恐惧。图 \ref{fig:tri} 展示了三级模型的一组时间序列和 $(k_1,k_2)$ 扫描结果。该扩展说明，多营养级恐惧效应可能进一步改变食物链稳定性和能量传递。

\begin{figure}[H]
    \centering
    \includegraphics[width=0.92\textwidth]{fig12_tritrophic_extension.png}
    \caption{三级食物链扩展模型的时间序列与参数扫描}
    \label{fig:tri}
\end{figure}

\setcounter{section}{8}
\section*{\centerline{八、结论}}

本文在经典捕食者--猎物模型基础上引入恐惧效应和记忆变量，构建了 M0 无恐惧、M1 即时恐惧和 M2 记忆恐惧三个层级模型，并实现了 M3 三级食物链扩展。理论分析表明，正共存平衡的存在依赖于捕食者维持条件
\[
\eta p>d_3h
\]
以及猎物有效增长条件
\[
r>d_1+d_2x^\ast.
\]
记忆变量不改变平衡点位置，因为平衡时 $z^\ast=y^\ast$，但它会改变线性化系统特征值，从而影响稳定性。

数值模拟显示，恐惧强度增加通常会降低捕食者平均密度和峰值；在部分参数区间内，高恐惧可使系统从周期振荡转向稳定共存；但在慢记忆区域，恐惧可能降低捕食者峰值却仍保留振荡。长记忆意味着风险反馈滞后，可能拉长或放大周期，甚至引入低密度风险。理论边界与数值热力图在主要稳定区域上基本一致，但在长周期和低密度风险附近存在差异。恐惧函数替换为指数型或二次有理型后，主要趋势仍保持，说明结论具有一定机制稳健性。

\setcounter{section}{9}
\section*{\centerline{九、问题与改进方向}}

本文仍有以下局限：
\begin{enumerate}[label=(\arabic*)]
    \item 没有进行真实参数拟合，所有结论都属于机制模拟和定性分析；
    \item 记忆变量是对空间记忆和行为记忆的 ODE 化简，不能直接替代 Yang 与 Jin 的 PDE 空间记忆模型；
    \item 当前分类依赖有限时间数值模拟，极长周期区域可能需要更长时间积分；
    \item 三级食物链只做扩展模拟，没有完整推导平衡点稳定性和分岔条件；
    \item 后续可加入 DDE 时滞、空间扩散或真实生态时间序列校准，但这些不应全部塞入主模型。
\end{enumerate}

\bibliographystyle{ieeetr}
\bibliography{refer}

\end{document}
